% Figure for Thermal and Statistical Physics §C4b.6 (Planck x^3/(e^x - 1) against Rayleigh-Jeans x^2, x = hbar omega/k_B T). Generated by th-C/tools/figs/c4b.py.
\begin{tikzpicture}[font=\small]
  \begin{axis}[nfig, width=8.6cm, height=5.6cm, xmin=0, xmax=10, ymin=0, ymax=3.2, xlabel={$x = \hbar\omega/k_BT$}, ylabel={scaled $u_\omega$}, legend style={font=\scriptsize, draw=none, fill=none, at={(0.98,0.97)}, anchor=north east}, legend cell align=left]
    \addplot[figB] coordinates {(0.0000,0.0000) (0.0833,0.0067) (0.1667,0.0255) (0.2500,0.0550) (0.3333,0.0936) (0.4167,0.1399) (0.5000,0.1927) (0.5833,0.2506) (0.6667,0.3126) (0.7500,0.3777) (0.8333,0.4448) (0.9167,0.5132) (1.0000,0.5820) (1.0833,0.6505) (1.1667,0.7181) (1.2500,0.7843) (1.3333,0.8485) (1.4167,0.9103) (1.5000,0.9694) (1.5833,1.0254) (1.6667,1.0780) (1.7500,1.1272) (1.8333,1.1727) (1.9167,1.2143) (2.0000,1.2521) (2.0833,1.2860) (2.1667,1.3160) (2.2500,1.3420) (2.3333,1.3642) (2.4167,1.3826) (2.5000,1.3973) (2.5833,1.4084) (2.6667,1.4160) (2.7500,1.4203) (2.8333,1.4214) (2.9167,1.4195) (3.0000,1.4147) (3.0833,1.4072) (3.1667,1.3971) (3.2500,1.3847) (3.3333,1.3701) (3.4167,1.3535) (3.5000,1.3350) (3.5833,1.3149) (3.6667,1.2931) (3.7500,1.2701) (3.8333,1.2458) (3.9167,1.2204) (4.0000,1.1941) (4.0833,1.1670) (4.1667,1.1392) (4.2500,1.1108) (4.3333,1.0821) (4.4167,1.0530) (4.5000,1.0237) (4.5833,0.9942) (4.6667,0.9648) (4.7500,0.9353) (4.8333,0.9060) (4.9167,0.8768) (5.0000,0.8480) (5.0833,0.8194) (5.1667,0.7912) (5.2500,0.7633) (5.3333,0.7360) (5.4167,0.7091) (5.5000,0.6827) (5.5833,0.6569) (5.6667,0.6317) (5.7500,0.6070) (5.8333,0.5830) (5.9167,0.5595) (6.0000,0.5367) (6.0833,0.5146) (6.1667,0.4931) (6.2500,0.4722) (6.3333,0.4520) (6.4167,0.4324) (6.5000,0.4135) (6.5833,0.3952) (6.6667,0.3776) (6.7500,0.3605) (6.8333,0.3441) (6.9167,0.3283) (7.0000,0.3131) (7.0833,0.2984) (7.1667,0.2843) (7.2500,0.2708) (7.3333,0.2578) (7.4167,0.2454) (7.5000,0.2335) (7.5833,0.2220) (7.6667,0.2111) (7.7500,0.2006) (7.8333,0.1906) (7.9167,0.1810) (8.0000,0.1718) (8.0833,0.1631) (8.1667,0.1547) (8.2500,0.1467) (8.3333,0.1391) (8.4167,0.1319) (8.5000,0.1250) (8.5833,0.1184) (8.6667,0.1121) (8.7500,0.1062) (8.8333,0.1005) (8.9167,0.0951) (9.0000,0.0900) (9.0833,0.0851) (9.1667,0.0805) (9.2500,0.0761) (9.3333,0.0719) (9.4167,0.0679) (9.5000,0.0642) (9.5833,0.0606) (9.6667,0.0572) (9.7500,0.0540) (9.8333,0.0510) (9.9167,0.0481) (10.0000,0.0454)};
    \addplot[figR, dashed] coordinates {(0.0000,0.0000) (0.0500,0.0025) (0.1000,0.0100) (0.1500,0.0225) (0.2000,0.0400) (0.2500,0.0625) (0.3000,0.0900) (0.3500,0.1225) (0.4000,0.1600) (0.4500,0.2025) (0.5000,0.2500) (0.5500,0.3025) (0.6000,0.3600) (0.6500,0.4225) (0.7000,0.4900) (0.7500,0.5625) (0.8000,0.6400) (0.8500,0.7225) (0.9000,0.8100) (0.9500,0.9025) (1.0000,1.0000) (1.0500,1.1025) (1.1000,1.2100) (1.1500,1.3225) (1.2000,1.4400) (1.2500,1.5625) (1.3000,1.6900) (1.3500,1.8225) (1.4000,1.9600) (1.4500,2.1025) (1.5000,2.2500) (1.5500,2.4025) (1.6000,2.5600) (1.6500,2.7225) (1.7000,2.8900) (1.7500,3.0625) (1.8000,3.2400) (1.8500,3.4225) (1.9000,3.6100) (1.9500,3.8025) (2.0000,4.0000)};
    \legend{Planck $x^3/(e^x - 1)$, Rayleigh--Jeans $x^2$}
    \node[font=\scriptsize, black!60, anchor=west] at (axis cs:3.1,1.62) {peak $x \simeq 2.82$};
  \end{axis}
\end{tikzpicture}
